
  \chapter{Antecedentes}
  \label{cap:antecedentes}


  \section{Complejidad cognitiva}
  \label{sec:cc-antecedentes}

  La medición de la complejidad del software tiene un antecedente clásico en la
  \emph{complejidad ciclomática} de McCabe~\cite{10.1109/TSE.1976.233837}, que
  cuenta el número de caminos linealmente independientes del grafo de control de
  un método. Se trata de una métrica útil para estimar el esfuerzo de prueba, pero
  no fue concebida para medir lo difícil que resulta comprender el código. En particular, trata de forma parecida estructuras que un lector humano no interpreta del mismo modo
  y no penaliza de forma explícita el anidamiento, que suele
  ser uno de los factores más costosos de seguir mentalmente.

  Para cubrir esa carencia, SonarSource propuso la \emph{complejidad
  cognitiva}~\cite{10.1145/3194164.3194186}, una métrica diseñada
  específicamente para aproximar el esfuerzo de comprensión. Su modelo se apoya
  en tres tipos de incremento: un incremento estructural por cada ruptura del flujo
  lineal de lectura, un incremento por anidamiento que penaliza más las
  estructuras cuanto más profundas están y un incremento asociado a ciertas
  secuencias de operadores lógicos. Es precisamente este componente de
  anidamiento el que convierte a los condicionales anidados en uno de los
  principales contribuyentes a la complejidad cognitiva y el que motiva el patrón
  objetivo de este trabajo.

  El uso de la complejidad cognitiva como aproximación a la comprensibilidad
  cuenta además con respaldo empírico. Muñoz Barón et
  al.~\cite{10.1145/3382494.3410636} encuentran, sobre un conjunto amplio de
  datos de comprensión, una correlación significativa entre esta métrica y el
  tiempo que tardan los desarrolladores en entender el código. Esto respalda su
  uso como criterio de mejora en RQ2.

  Conviene matizar, no obstante, que la métrica empleada en este trabajo no es la
  implementación oficial de SonarQube, sino una implementación \emph{proxy} del
  modelo (\autoref{sec:proxy}). Su validez de constructo se acota mediante el
  contraste descrito en la \autoref{sec:sonar-val}.

  \section{Refactorización y preservación semántica}
  \label{sec:refactor-antecedentes}

  Se entiende por refactorización el conjunto de transformaciones que mejoran
  la estructura interna del código sin alterar su comportamiento observable.
  La referencia clásica en este ámbito es Fowler~\cite{fowler2018refactoring},
  que organiza la disciplina en torno a un catálogo de transformaciones pequeñas
  y con nombre propio, aplicadas de forma incremental. Entre ellas se encuentra
  la consolidación de expresiones condicionales, que es precisamente la familia a
  la que pertenece el patrón estudiado en este trabajo: fusionar condicionales
  anidados en una única condición conjuntiva.

  La idea de que una refactorización debe preservar el comportamiento se remonta
  a la tesis de Opdyke~\cite{opdyke1992refactoring}, donde estas transformaciones
  se definen como operaciones que conservan el comportamiento del programa
  \emph{siempre que se cumplan sus precondiciones}. Este principio es el que
  guía el enfoque del prototipo. En lugar de asumir que una transformación es
  segura en abstracto, su aplicación se restringe a aquellos casos en los que un
  conjunto de precondiciones verificables, las condiciones P1--P5 de la
  \autoref{sec:precondiciones}, permite garantizar la equivalencia.

  La noción de equivalencia que se adopta en este trabajo es la de equivalencia
  observable: el código refactorizado debe producir el mismo efecto que el
  original para cualquier entrada. En el caso del patrón
  \lstinline|if(A){if(B)}| frente a \lstinline|if(A && B)|, esta equivalencia
  depende de forma crítica de que las condiciones no introduzcan efectos
  colaterales ni alteren la semántica de cortocircuito.

  El trabajo de Mens y Tourwé~\cite{10.1109/TSE.2004.1265817} ofrece una visión
  general de la disciplina y distingue varias actividades dentro del proceso de
  refactorización: identificar dónde refactorizar, comprobar que la
  transformación preserva el comportamiento, aplicarla y evaluar su efecto sobre
  la calidad del software. Esta descomposición encaja con la estructura del
  prototipo desarrollado en este TFM, que se organiza precisamente en torno a
  detección, verificación de precondiciones, transformación y medición del
  impacto.
  La
  \autoref{tab:actividades-pipeline} hace explícita esta correspondencia.

  \begin{table}[ht]
  \centering
  \caption{Correspondencia entre las actividades de refactorización descritas
  por Mens y Tourwé~\cite{10.1109/TSE.2004.1265817} y las etapas del prototipo.}
  \label{tab:actividades-pipeline}
  \begin{tabular}{@{}lll@{}}
  \toprule
  \textbf{Actividad (Mens y Tourwé)} & \textbf{Etapa del prototipo} & \textbf{Componente} \\
  \midrule
  Identificar dónde refactorizar   & Detección del patrón      & \texttt{NestedIfDetector} \\
  Garantizar la preservación       & Verificación de P1--P5    & Precondiciones \\
  Aplicar la transformación        & Refactorización           & \texttt{NestedIfTransformer} \\
  Evaluar el efecto                & Medición del impacto      & CC \emph{proxy} \\
  \bottomrule
  \end{tabular}
  \end{table}

  \section{Análisis estático sobre AST con JavaParser}
  \label{sec:ast-antecedentes}

  Detectar y transformar el patrón objetivo exige operar sobre la estructura
  sintáctica del código, no sobre su representación textual. La forma habitual de abordar este tipo de problema es mediante un árbol de sintaxis abstracta (AST)
  , en el que cada construcción del lenguaje, por ejemplo, una sentencia
  \lstinline|if|, una expresión binaria o un bloque, se representa como un
  nodo con una posición bien definida dentro de la jerarquía del programa.
  Trabajar sobre el AST evita la fragilidad de los enfoques basados en texto y
  permite razonar de forma directa sobre relaciones estructurales como el
  anidamiento. La \autoref{fig:ast-anidado} ilustra esta idea sobre el patrón
  objetivo: el anidamiento sintáctico de los dos \lstinline|if| se traduce en una
  relación de descendencia entre nodos del árbol, de modo que el \lstinline|if|
  interno aparece como descendiente de la rama \lstinline|then| del externo.

  \begin{figure}[ht]
  \centering
  \includegraphics[width=0.45\linewidth]{figuras/img/AST.png}
  \caption{Árbol de sintaxis abstracta (AST) del patrón objetivo de dos
  sentencias \lstinline|if| anidadas. El anidamiento se representa como una
  relación de descendencia: el \lstinline|if| interno cuelga de la rama
  \lstinline|then| del externo, y el cuerpo \lstinline|S| queda en el nivel más
  profundo.}
  \label{fig:ast-anidado}
  \end{figure}

  El prototipo desarrollado en este trabajo se apoya en
  JavaParser~\cite{smith2023javaparser}, una biblioteca que permite construir el
  AST de código fuente Java y recorrerlo o modificarlo programáticamente. La
  versión concreta utilizada se fija en \texttt{pom.xml}
  (\autoref{sec:experimento}). Sobre esta base, la fase de detección recorre el
  árbol en busca de sentencias \lstinline|if| anidadas que cumplan el patrón
  objetivo, mientras que la fase de transformación reescribe el árbol
  \emph{in situ}: combina ambas condiciones en una conjunción y reubica el cuerpo
  bajo el nuevo nodo. Esta manipulación se mantiene deliberadamente
  conservadora: no reordena expresiones ni intenta simplificarlas, y preserva la
  semántica de precedencia envolviendo en paréntesis las condiciones cuando es
  necesario, por ejemplo cuando una disyunción pasa a formar parte de una
  conjunción.

  La limitación relevante para este trabajo es que el análisis realizado es
  fundamentalmente sintáctico. El AST describe la forma del programa, pero no
  resuelve por sí mismo información semántica como los tipos, los alias o el
  flujo de datos. En consecuencia, a partir del árbol no siempre puede decidirse
  si una llamada a método introduce efectos colaterales o si dos expresiones son
  equivalentes desde el punto de vista observacional. Esta frontera justifica la
  estrategia conservadora del prototipo: en lugar de intentar un análisis
  semántico completo, la aplicabilidad se restringe mediante precondiciones
  comprobables sobre la estructura y una lista explícita de construcciones
  admitidas (\autoref{sec:precondiciones}), descartando todo aquello que no pueda
  verificarse con garantías.

  \section{Grandes modelos de lenguaje aplicados a tareas de código}
  \label{sec:llm-antecedentes}

  La aplicación de grandes modelos de lenguaje (LLM) a tareas de programación
  se consolidó con Codex~\cite{chen2021codex}, un modelo entrenado sobre código
  público que demostró una capacidad notable para sintetizar programas a partir
  de descripciones en lenguaje natural y que dio lugar a asistentes de
  programación de uso masivo. Desde entonces, el interés se ha desplazado desde
  la generación de código nuevo hacia tareas sobre código existente, entre ellas
  la refactorización, que es precisamente el escenario abordado en RQ3.

  La evidencia empírica sobre refactorización con LLM es prometedora, aunque
  todavía debe interpretarse con cautela. DePalma
  et al.~\cite{10.1016/j.eswa.2024.123602} observan, sobre un conjunto de
  fragmentos de código Java, que ChatGPT propone mejoras útiles en la mayoría de
  los casos, si bien con resultados desiguales según el atributo de calidad
  considerado y sin una garantía sistemática de preservación del comportamiento.
  Liu et al.~\cite{liu2024llmrefactoring}, por su parte, encuentran que los
  modelos identifican por sí solos una fracción reducida de las oportunidades
  reales de refactorización y que su rendimiento depende en gran medida del modo
  en que la tarea se formule y acote mediante el \emph{prompt}. En la misma línea,
  Cordeiro et al.~\cite{cordeiro2024llmrefactoring} estudian la capacidad de
  refactorización de los LLMs sobre proyectos Apache reales y observan que,
  aunque los modelos producen código con una estructura mejorada, no garantizan la
  equivalencia semántica de la transformación.

  De este panorama se derivan dos implicaciones que condicionan el diseño de
  RQ3. La primera es que la corrección de una refactorización generada por un
  LLM no puede darse por supuesta, sino que debe comprobarse frente a una
  referencia fiable. En este trabajo, esa referencia la proporcionan el
  \emph{baseline} determinista del prototipo y el oráculo automático
  (\autoref{sec:oraculo}). La segunda es que la salida de estos modelos puede variar entre ejecuciones, lo que obliga a fijar un protocolo reproducible, modelo, temperatura, versión del \emph{prompt} y número de intentos por caso; para que la evaluación resulte interpretable
  (\autoref{sec:exp-rq3}).

  \section{Síntesis crítica orientada al problema}
  \label{sec:sintesis-antecedentes}

  Los antecedentes revisados convergen en el problema que aborda este trabajo,
  pero ninguno lo cubre de forma completa. La complejidad cognitiva
  (\autoref{sec:cc-antecedentes}) aporta la métrica y su justificación
  empírica; la teoría de la refactorización
  (\autoref{sec:refactor-antecedentes}) proporciona el principio de
  preservación del comportamiento condicionado al cumplimiento de
  precondiciones; el análisis sobre AST
  (\autoref{sec:ast-antecedentes}) ofrece la base tecnológica, aunque dentro
  de una frontera esencialmente sintáctica; y la literatura sobre LLM
  (\autoref{sec:llm-antecedentes}) delimita tanto el potencial de estos
  modelos como la necesidad de validar sus salidas antes de aceptarlas como
  refactorizaciones correctas.

  El trabajo más cercano al presente TFM es el de Saborido et
  al.~\cite{10.1109/ACCESS.2022.3144743}, centrado en la reducción automática
  de la complejidad cognitiva. Su propuesta modela el problema como una
  búsqueda sobre secuencias de extracción de método (\emph{Extract Method}) y
  la evalúa sobre diez proyectos de código abierto, logrando resolver una
  fracción elevada de las incidencias de complejidad cognitiva detectadas por
  SonarQube. El presente trabajo comparte con ese enfoque el objetivo general
  de reducir la complejidad cognitiva mediante refactorización automática y se
  sitúa en el mismo contexto experimental de proyectos open source
  (\autoref{sec:corpus}). Sin embargo, difiere tanto en el tipo de
  transformación como en el alcance. Mientras que Saborido et al.
  \emph{descomponen} métodos mediante extracción, este trabajo
  \emph{combina} condicionales anidados en una única condición. Se trata, por
  tanto, de operaciones distintas pero complementarias sobre una misma
  métrica. Además, aquel trabajo no aborda la capacidad de los modelos de
  lenguaje para realizar la transformación, que es precisamente el objeto de
  RQ3.

  Conviene situar también el trabajo frente a herramientas de refactorización
  automática de uso extendido. Los entornos de desarrollo integrados (IntelliJ
  IDEA, Eclipse) ofrecen un repertorio amplio de refactorizaciones, entre ellas
  la consolidación de condicionales, pero las accionan el desarrollador
  manualmente y no utilizan la complejidad cognitiva como criterio de
  decisión. OpenRewrite es un \emph{framework} de transformación AST que
  permite definir y ejecutar \emph{recetas} automáticas con garantías formales por
  receta; sin embargo, no incluye, en su catálogo estándar, una receta orientada
  específicamente a la combinación de condicionales anidados con criterio de CC y
  comprobación de precondiciones. Las reglas automáticas de SonarQube
  (\emph{quick fixes}) se limitan a sugerencias de baja complejidad y no cubren
  este patrón. Ninguna de estas herramientas aúna los tres rasgos del presente
  trabajo: transformación formalizada con precondiciones verificables, medición
  de impacto en CC y evaluación de los LLM como agente alternativo.

  La \autoref{tab:comparativa} resume este posicionamiento frente a enfoques
  previos. El hueco que cubre este trabajo se sitúa en la intersección de tres
  rasgos que, por separado, sí aparecen en la literatura: una transformación
  concreta acotada mediante precondiciones \emph{verificables}, el uso de la
  complejidad cognitiva como criterio de impacto y una evaluación explícita de
  la capacidad de los LLM para reproducir la refactorización frente a un
  \emph{baseline} determinista.

  \begin{table}[H]
  \centering
  \caption{Posicionamiento del trabajo frente a enfoques previos de refactorización y reducción de complejidad.}
  \label{tab:comparativa}
  \begin{tabular}{@{}p{3.1cm}p{2.8cm}p{2.8cm}p{2.5cm}@{}}
  \toprule
  \textbf{Enfoque} & \textbf{Transformación} & \textbf{Garantía de seguridad} & \textbf{Papel de los LLM} \\
  \midrule
  Refactors de IDE / catálogo de Fowler~\cite{fowler2018refactoring}
    & Repertorio amplio
    & Heurística de la herramienta
    & No considerado \\
  Saborido et al.~\cite{10.1109/ACCESS.2022.3144743}
    & Extracción de método (búsqueda)
    & Refactor que preserva el comportamiento
    & No considerado \\
  LLM para refactor~\cite{10.1016/j.eswa.2024.123602, liu2024llmrefactoring}
    & Abierta
    & Sin garantía formal
    & Constituyen el objeto de análisis \\
  \textbf{Este trabajo}
    & Combinación de condicionales anidados
    & Precondiciones P1--P5 verificables
    & Evaluados en RQ3 \\
  \bottomrule
  \end{tabular}
  \end{table}
